% TOP_PROBLEM: {"id": 3000079, "title": "Small quasi-kernels in directed graphs", "category_id": 2, "category": {"id": 2, "name": "combinatorics", "display_name": "Combinatorics", "description": "Counting problems, graph theory, discrete structures.", "slug": "combinatorics", "order_index": 2, "created_at": "2026-07-31T15:26:25.671Z"}, "status": "partially_solved"}
% FIRST_POSTED: null
% ENTRY_KIND: statement_only
% Imported from: batch09.tex; UnsolvedMath ID 3000079
\documentclass[11pt]{article}
\usepackage[margin=25mm]{geometry}
\usepackage{fontspec}
\setmainfont{Latin Modern Roman}
\usepackage{amsmath,amssymb,mathtools}
\usepackage{unicode-math}
\setmathfont{Latin Modern Math}
\usepackage{xurl,hyperref}
\hypersetup{colorlinks=true,urlcolor=blue}
\setcounter{secnumdepth}{0}
\setlength{\parindent}{0pt}
\setlength{\parskip}{5pt}
\setlength{\emergencystretch}{3em}
\newcommand{\sref}[2]{\hyperlink{source:#1:#2}{[S#2]}}
\newcommand{\cataloguescope}[1]{\paragraph{Recorded target.}#1}
\begin{document}
% UnsolvedMath ID 3000079; source catalogue code AMR-029-0079
\setcounter{section}{3000078}
\section{Small quasi-kernels in directed graphs}\label{3000079}
{\small\sffamily UnsolvedMath ID 3000079 · Combinatorics}
\subsection{Definitions and mathematical statement}

\paragraph{Shared notation.}$\mathbb N=\{1,2,\ldots\}$, and $\mathbb Z,\mathbb Q,\mathbb R,\mathbb C$ denote the integers, rationals, reals and complex numbers. Any other conventions are specified locally.


Let $D=(V,A)$ be a finite loopless digraph with $n=|V|$ and $d^+(v)\ge1$ for every vertex. Here a quasi-kernel is an independent set $Q\subseteq V$ such that every vertex $v$ has a directed path of length at most two from $v$ to some $q\in Q$. Independence means there is no arc in either direction between two members of $Q$. Length-zero paths are allowed; thus vertices in $Q$ satisfy the condition themselves. This is the inward convention used with the positive-outdegree hypothesis.
\paragraph{Statement recorded in the supplied source.}
Does every such sink-free digraph have a quasi-kernel $Q$ with $|Q|\le n/2$? \sref{3000079}{2}
\subsection{Short English statement}
Does every finite directed graph without sinks have an independent set of at most half its vertices reachable from every vertex within two steps?
\subsection{Sources}
\begin{description}
\item[\hypertarget{source:3000079:1}{S1}] ULAM.AI and the UnsolvedMath Contributors, \emph{UnsolvedMath: A Curated Collection of Open Mathematics Problems}, record 3000079 (AMR-029-0079). CC BY 4.0. \url{https://huggingface.co/datasets/ulamai/UnsolvedMath}.
\item[\hypertarget{source:3000079:2}{S2}] H. Langlois, F. Meunier, R. Rizzi, S. Vialette and Y. Zhou, Quasi-kernels in split graphs (2024), Introduction, quasi-kernel definition and Conjecture 1, p.1. \url{https://arxiv.org/pdf/2312.15519}.
\end{description}
\label{end:3000079}
\end{document}
